% !TEX TS-program = lualatex+se
% !TEX encoding = UTF-8

\documentclass[oneside, 12pt, a4paper]{book}
%\usepackage[letterpaper,showframe]{geometry}
\usepackage{fullpage} % redukcja marginesów
\usepackage{geometry}
\geometry{
	left=10mm,
	right=10mm,
	top=12mm,
	bottom=12mm,
	marginparwidth=0mm,
	marginparsep=0mm
	}
% --- GREGORIO ---	
\usepackage{gregoriotex}
\gresetlinecolor{gregoriocolor}
\gresetheadercapture{annotation}{greannotation}{string}% odczytuje adnotacje nad inicjałem jako ciąg, co pozwala na jego renderowanie
\gresetheadercapture{commentary}{grecommentary}{string}% odczytuje komentarz jako ciąg, co pozwala na jego renderowanie
\grechangedim{commentaryraise}{0.3cm}{scalable} % podnosi komentarz wystarczająco ponad górną linię
\grechangestyle{annotation}{\color{black}\small}% format adnotacji
\grechangestyle{commentary}{\color{gregoriocolor}\small}% format komentarza
\grechangestyle{abovelinestext}{\color{gregoriocolor}\small}% format napisu nad liniami
\grechangestyle{initial}{\fontsize{36}{36}\selectfont} % zmniejsza wielkość inicjału do bardziej proporcjonalnego rozmiaru (=36pt)
\grechangedim{intersyllablespacenotes}{0.26cm}{scalable}% Minimalna odległość między nutami różnych sylab.

% umieszczenie zawartości składu ich unikalnego argumentu po lewej i prawej stronie każdej linii
\let\grelocalleftbox\localleftbox
\let\grelocalrightbox\localrightbox

%\newcommand{\separator}{\greseparator{3}{40}}
\newfontface\vectis{VectisMonumentalW01-Regular}
\def\greinitialformat#1{\fontsize{20}{20}\vectis #1}
% --- format pierwszej linii 
% \setfirstannotation{\fontsize{10}{10}\textcolor{red}{Ant.}}%
% --- format drugiej linii
% \setsecondannotation{\fontsize{10}{10}\sc{\textcolor{red}8.}}%
% --- dodawanie spacji
% \tolerance=9999 %(or less)
% \pretolerance=500  %default is -1  This will spread your chant out over the bars
% --- wymuś hypheny między sylabami
% \grechangedim{\gremaximumspacewithoutdash}{0pt}
% --- redukcja spacji między sylabami i glyphami: \hspace{-.2ex}
% --- lub wstam sylaby lub glyphy między to: \mbox{...}
% --- lub użyj ( z), aby przesunąć neum bliżej początku pięciolinii. Im więcej spacji w nawiasach, tym bardziej zawęża spację.
% ------------------

%  Konfiguracja kolumn
\usepackage{paracol}
\setcolumnwidth{13cm,5cm}
\setlength{\columnsep}{0.5cm}
\footnotelayout{m}
% \setlength{\columnseprule}{0.4pt}
% \colseprulecolor{red}
\usepackage{tcolorbox}
\newtcolorbox{nutybox}{size=minimal, colback=white, colframe=red, arc=0mm, width=\linewidth, halign=left, valign=center}
\newtcolorbox{tlumbox}{size=minimal, colback=white, colframe=red, arc=0mm, width=\linewidth, halign=left, valign=top, fontupper=\linespread{.8}\selectfont}

\usepackage{Tabbing}
\usepackage{longtable}
\usepackage[savepos]{zref}

\usepackage[latin,polish]{babel}
\usepackage{fontspec}
%\usepackage{times}
%\usepackage[math]{anttor}
%\usepackage{tgtermes}
%\usepackage{mathptmx}
\usepackage{ebgaramond}
\setmainfont{EB Garamond}
%\usepackage{libertine}
%\setmainfont{Libertinus Serif}
\defaultfontfeatures{Ligatures=TeX}

\pagestyle{plain}	% usuwanie nagłówków/stopek z\,dokumentu
\usepackage[indent=0pt]{parskip} % ustawianie paragrafów
\usepackage{setspace} % odstępy między liniami
%\setstretch{1.5} % półtorej linii odległości między wierszami

\usepackage{graphicx} % umożliwia import obrazów
\graphicspath{{img/}} % katalog obrazów
\usepackage{pgfornament}% ornamenty graficzne

\usepackage{titlesec} % ustawianie nagłówków

% --- komendy własne ---
\newcommand{\red}[1]{\textcolor{red}{#1}}
\newcommand{\black}[1]{\textcolor{black}{#1}}

\newcommand{\koniec}{
	\begin{center}
%		\vspace{5mm}
%		\pgfornament[width=0.6cm]{#1}
		\newpage
	\end{center}}

\begin{document}
\titlespacing*{\chapter}{0pt}{0pt}{5pt}	
\titleformat{\chapter}[hang]{\filcenter\LARGE\scshape\bfseries}{}{0pt}{} % format nagłówków
\titlespacing*{\section}{0pt}{0pt}{0pt}	
\titleformat{\section}[hang]{\filcenter\large\scshape}{}{0pt}{} % format nagłówków

\setlength{\parindent}{0pt}
\gresetbracerendering{font}

\pagenumbering{gobble} % nie jest wyprowadzany żaden numer strony, chociaż numer jest nadal resetowany. Odniesienia do tej strony również są puste.
\large
\begin{spacing}{0.8}
\begin{center}
	\section*{Veni veni Emmanuel}
	\small\textit{Psalteriolum Cantionum Catholicarum, Kolonia, 1710r.}
\end{center}	
\end{spacing}
\begin{paracol}{2}
	\ensurevspace{4cm}
	\begin{nutybox}
	\greannotation{Hymnus}\greannotation{1.}\gabcsnippet{(cb4) VE(d)ni,(f') ve(h)ni,(h) Em(h)má(gih)nu(g)el,(f_) (,) Cap(g)tí(h')vum(f) sol(d)ve(f) Is(gvED)ra(c)el,(d.) (;) Qui(g) ge(g')mit(d) in(d) ex(e)sí(f_e)li(d)o(c_) (,) Pri(f)vá(g')tus(h) De(h)i(h) Fí(gih)li(g)o.(f.) (::) <c><sp>R/</sp>.</c> Gau(j)de,(j_) gau(h)de,(h_) Em(h)má(gih)nu(g)el(f_) (,) Na(g)scé(h')tur(f) pro(d) te,(f) Is(gvED)ra(c)el.(d.) (::)}
	\end{nutybox}
	\switchcolumn
	\begin{tlumbox}
	\vspace{3.5mm} 
	\red{1}.~O przyjdź, o przyjdź, Emanuelu,
z niewoli wyrwij Izraela;
za grzech wygnanie cierpi lud.
Jednorodzony Synu Boży!
\red{\Rbar.}~Rozraduj się, rozraduj się,
bo przyjdzie twój Emmanuel!
	\end{tlumbox}
  	\ensurevspace{4cm}
	\switchcolumn*
	\begin{nutybox}
	\gresetinitiallines{0}\gabcsnippet{(cb4)  <c>2.</c> Ve(d)ni,(f') ve(h)ni,(h) Rex(h) gén(gih)ti(g)um,(f_) (,) Ve(g)ni,(h') Red(f)ém(d)ptor(f) óm(gvED)ni(c)um:(d.) (;) Ut(g) sal(g')ves(d) tu(d)os(e) fá(f_e)mu(d)los(c_) (,) Pec(f)cá(g')ti(h) si(h)bi(h) cón(gih)sci(g)os.(f.) (::) <c><sp>R/</sp>.</c> Gau(j)de,(j_) gau(h)de...(h_)}
	\end{nutybox}
	\switchcolumn
	\begin{tlumbox}
	\vspace{3.5mm} 
	\red{2}.~O przyjdź, o przyjdź, o Królu królów,
przyjdź Zbawicielu, przyjdź i ocal świat,
uwolnij nas z niewoli zła 
i Twoją bliskość odczuć nam daj.
	\end{tlumbox}
	\ensurevspace{4cm}
	\switchcolumn*
	\begin{nutybox}
	\gresetinitiallines{0}\gabcsnippet{(cb4)  <c>3.</c> Ve(d)ni,(f') ve(h)ni,(h) O(h) O(gih)ri(g)ens,(f_) (,) So(g)lá(h')re(f) nos(d) ad(f)vé(gvED)ni(c)ens:(d.) (;) No(g)ctis(g') de(d)pél(d)le(e) né(f_e)bu(d)las,(c_) (,) Di(f)rás(g')que(h) no(h)ctis(h) té(gih)ne(g)bras.(f.) (::) <c><sp>R/</sp>.</c> Gau(j)de,(j_) gau(h)de...(h_)}
	\end{nutybox}
	\switchcolumn
	\begin{tlumbox}
	\vspace{3.5mm} 
	\red{3}.~O przyjdź, o przyjdź, o Wschodzie Słońca,
i oświeć ziemię Światłem Twym,
niech przezwycięży nocne mgły
rozproszy ciemność nocy mej.
	\end{tlumbox}
	\ensurevspace{4cm}
	\switchcolumn*
	\begin{nutybox}
	\gresetinitiallines{0}\gabcsnippet{(cb4)  <c>4.</c> Ve(d)ni,(f') Cla(h)vis(h) Da(h)ví(gih)di(g)ca,(f_) (,) Re(g)gna(h') re(f)clú(d)de(f) c<sp>'ae</sp>(gvED)li(c)ca;(d.) (;) Fac(g) i(g')ter(d) tu(d)tum(e) sú(f_e)pe(d)rum,(c_) (,) Et(f) clau(g')de(h) vi(h)as(h) ín(gih)fe(g)rum.(f.) (::) <c><sp>R/</sp>.</c> Gau(j)de,(j_) gau(h)de...(h_)}
	\end{nutybox}
	\switchcolumn
	\begin{tlumbox}
	\vspace{3.5mm}
	\red{4}.~O przyjdź, o przyjdź, Kluczu Dawida,
niebo otworzyć w mocy Twej jest;
bezpieczną drogą prowadź wzwyż
na zawsze zamknij piekła drzwi.
	\end{tlumbox}
	\ensurevspace{4cm}
	\switchcolumn*
	\begin{nutybox}
	\gresetinitiallines{0}\gabcsnippet{(cb4)  <c>5.</c> Ve(d)ni,(f') O(h) Jes(h)se(h) Vír(gih)gu(g)la:(f_) (,) Ex(g) ho(h')stis(f) tu(d)os(f) ún(gvED)gu(c)la,(d.) (;) De(g) spe(g')cu(d) tu(d)os(e) tár(f_e)ta(d)ri(c_) (,) E(f)duc,(g') et(h) an(h)tro(h) bá(gih)ra(g)thri.(f.) (::) <c><sp>R/</sp>.</c> Gau(j)de,(j_) gau(h)de...(h_)}
	\end{nutybox}
	\switchcolumn
	\begin{tlumbox}
	\vspace{3.5mm}
	\red{5}.~O przyjdź, o przyjdź, Różdżko Jessego,
z rąk nieprzyjaciół wyzwól nas,
z przepaści i otchłani zła
z czeluści śmierci wyrwij nas.
%\footnotemark
	\end{tlumbox}	
	\ensurevspace{4cm}
	\switchcolumn*
	\begin{nutybox}
	\gresetinitiallines{0}\gabcsnippet{(cb4)  <c>6.</c> Ve(d)ni,(f') ve(h)ni,(h) Ad(h)o(gih)ná(g)i,(f_) (,) Qui(g) pó(h')pu(f)lo(d) in(f) Sí(gvED)na(c)i,(d.) (;) Le(g)gem(g') de(d)dí(d)sti(e) vér(f_e)ti(d)ce,(c_) (,) In(f) ma(g')je(h)stá(h)te(h) gló(gih)ri(g)æ.(f.) (::) <c><sp>R/</sp>.</c> Gau(j)de,(j_) gau(h)de,(h_) Em(h)má(gih)nu(g)el(f_) (,) Na(g)scé(h')tur(f) pro(d) te,(f) Is(gvED)ra(c)el.(d.) (::)}
	\end{nutybox}
	\switchcolumn
	\begin{tlumbox}
	\vspace{3.5mm}
	\red{6}.~O przyjdź, o przyjdź, o święty Adonaj,
coś Twemu ludowi Prawo dał
na górze Synaj, najwyższej z gór,
niech Słowo Twe na nowo stworzy nas.
\red{\Rbar.}~Rozraduj się, rozraduj się,
bo przyjdzie twój Emmanuel!
%\footnotemark
	\end{tlumbox}	
	\ensurevspace{4cm}
	\switchcolumn*
	\begin{nutybox}
	\gresetinitiallines{0}\gabcsnippet{(cb4)  <c>7.</c> Ve(d)ni,(f') O(h) Sa(h)pi(h)én(gih)ti(g)a,(f_) (,) Quæ(g) hic(h') di(f)spó(d)nis(f) óm(gvED)ni(c)a:(d.) (;) Ve(g)ni,(g') vi(d)am(d) pru(e)dén(f_e)ti(d)æ(c_) (,) Ut(f) dó(g')ce(h)as(h) et(h) gló(gih)ri(g)æ.(f.) (::) <c><sp>R/</sp>.</c> Gau(j)de,(j_) gau(h)de...(h_)}
	\end{nutybox}
	\switchcolumn
	\begin{tlumbox}
	\vspace{3.5mm}
	\red{7}.~O przyjdź, Mądrości, Słowo Boże,
nad wszystkim obejmij rządy Twe.
Naucz nas roztropności, przyjdź,
ścieżkami życia prowadź nas.\footnotemark
	\end{tlumbox}	
	\footnotetext{tłum. ze strony: https://www.pjthe3.com/o-przyjdz-o-przyjdz-emmanuelu/} 
\end{paracol}	
\end{document}